\begin{statement}{A.21}{Proved}{The diagonal elliptic Soul\'e character at the split prime three}
Let \(k=\mathbb Q(\sqrt{-2})\), let
\(3=\mathfrak p\overline{\mathfrak p}\), and let \(E/k\)
be a CM elliptic curve with endomorphism ring
\(\mathbb Z[\sqrt{-2}]\), good reduction at three, and
its standard puncture at the origin. In the homological
normalization \(\psi\overline\psi=\chi_3\), the diagonal
elliptic Soul\'e character
\(\kappa_{(2,2)}\) is nonzero after tensoring with
\(\mathbb Q_3\). Its restriction to any finite-index
subgroup of its defining Galois group remains nonzero.
This assertion at three is proved below; it is not an
application of a theorem whose statement assumes
\(p\ge5\).
\end{statement}
\begin{proof}
Write \(S\) for the two primes above three and the
infinite places. We first prove the precise cohomological
input, including the residual character:
\[
 H^2(G_{k,S},\mathbb Z_3(2))=0.                         \tag{A.21.1}
\]
Put \(L=k(\zeta_3)=\mathbb Q(\sqrt{-2},\sqrt{-3})\).
The discriminants of its three quadratic subfields are
\(-8,-3,24\); hence \(\operatorname{disc}(L)=576\).
Its Minkowski bound is \(36/\pi^2<4\). There is no
prime ideal of norm two: at two the extension supplied
by \(\sqrt{-3}\) is the unramified quadratic extension,
so the residue degree is two. The only prime ideals of
norm below four are the two primes above three. Each
has square principal, generated by
\(1+\sqrt{-2}\) or \(1-\sqrt{-2}\). Thus the class
group is generated by elements of order dividing two,
and \(\operatorname{Cl}(L)[3]=0\).

There are exactly two primes in \(S_L\), both fixed by
\(A=\operatorname{Gal}(L/k)\simeq C_2\). Kummer theory
and the Brauer invariant-sum sequence give
\[
 H^2(G_{L,S_L},\mu_3)
  =\ker(\mathbb F_3^2\xrightarrow{\sum}\mathbb F_3)
  \simeq\mathbb F_3,
\]
with trivial \(A\)-action. Passing from \(\mu_3\) to
\(\mathbb F_3(2)\) tensors this action by the nontrivial
cyclotomic character of \(A\). Averaging over \(A\)
is exact, and therefore
\(H^2(G_{k,S},\mathbb F_3(2))=0\). The coefficient
sequences prove vanishing for every
\(\mathbb Z/3^n(2)\). Finite-ramification cohomology
with finite coefficients is finite, so the
Mittag--Leffler property and continuous comparison
prove (A.21.1). This uses the same precise finiteness
and comparison theorems as D.1.

We next explain why the cyclotomic elliptic-unit map
has nonzero rational image, without any integral
claim that twelve is a three-adic unit. Put
\(L_\infty=k(\mu_{3^\infty})\),
\(\Delta=\operatorname{Gal}(k(\mu_3)/k)\simeq C_2\),
and \(\Gamma=\operatorname{Gal}(L_\infty/k(\mu_3))\simeq\mathbb Z_3\).
Let \(\mathcal E_L\) be the inverse limit under norms
of the three-adic completions of global units in this
tower, and let \(\mathcal A_L\) be the inverse limit
of the three-primary ideal class groups. A superscript
\((1)\) in the next formulas means the cyclotomic
Tate twist. The norm-and-Kummer construction gives
\[
 e_L:\mathcal E_L(1)_{\Delta\times\Gamma}
              \longrightarrow H^1(G_{k,S},\mathbb Z_3(2)).
                                                               \tag{A.21.2}
\]
Both its kernel and the twisted class-group coinvariants
are finite. Here is the cohomological verification of
the rational assertion needed, so that no restriction
\(p\ne3\) from a regulator theorem is imported.

Form the inverse corestriction complex
\(C_{\rm Iw}=R\varprojlim_nR\Gamma(G_{L_n,S},\mathbb Z_3(1))\).
Shapiro's cochain isomorphism, followed by the two-term
resolution for \(\Gamma\), identifies its derived
specialization at the character \(\chi_3^{-1}\), with
the corresponding \(\Delta\)-component, with
\(R\Gamma(G_{k,S},\mathbb Z_3(2))\).
The finite-level Kummer sequences identify its first
cohomology with norm-compatible \(S\)-units and its
second cohomology with an extension of
\(\mathcal A_L\) by the local Brauer-invariant modules.
The valuation and Brauer modules at the two primes
above three have trivial \(\Delta\)-action, whereas
\(\chi_3^{-1}|_\Delta\) is nontrivial. They therefore
vanish in this component; ordinary units and
\(S\)-units give the same component. The resulting
one-variable specialization sequences are
\[
\begin{split}
0&\longrightarrow\mathcal E_L(1)_{\Delta\times\Gamma}
 \longrightarrow H^1(G_{k,S},\mathbb Z_3(2))
 \longrightarrow\mathcal A_L(1)^{\Delta\times\Gamma}
 \longrightarrow0,\\
H^2(G_{k,S},\mathbb Z_3(2))
 &\simeq\mathcal A_L(1)_{\Delta\times\Gamma},
\end{split}                                                   \tag{A.21.3}
\]
up to finite groups, which is sufficient here. More
explicitly, any finite error from passing between
finite-level Kummer sequences and completed units
vanishes after tensoring with \(\mathbb Q_3\); the
completed cohomology groups are compact; inverse limits
of their exact sequences are exact by compactness. At
each finite coefficient level the groups are finite,
so the corresponding diagonal towers are Mittag--Leffler. The maps in (A.21.3) are the
norm-and-Kummer maps defining (A.21.2), not different
abstract comparison maps. The class-group Iwasawa
module is finitely generated torsion. For such a
one-variable module, finite twisted coinvariants imply
finite twisted invariants, by its elementary-divisor
decomposition. Equation (A.21.1) proves this finiteness
and makes (A.21.2) an isomorphism after tensoring with
\(\mathbb Q_3\). The global Euler characteristic formula
then gives
\(\dim_{\mathbb Q_3}H^1(G_{k,S},\mathbb Q_3(2))=1\).
Thus the rational image of (A.21.2) is nonzero.

Let \(\mathcal C_L^{\rm R}\subset\mathcal E_L\) be
Rubin's elliptic-unit module. Rubin's original
hypotheses hold at three: the Hilbert class field of
\(k\) is \(k\), its roots of unity have order two,
\([k(\mu_3):k]=2\), and three splits in \(k\).
In particular \cite[\S4, Theorem 4.1(i)]{DE_Rubin1991}
allows this prime; its restriction is
\(p\nmid\#\mu(H)[k_0:k]\), not \(p\nmid6\).
The first characteristic-ideal equality in that
theorem identifies
\(\operatorname{char}(\mathcal E_L/\mathcal C_L^{\rm R})\)
with \(\operatorname{char}(\mathcal A_L)\).
The finite specialization in (A.21.3), the
one-variable torsion structure theorem, and
\cite[Corollary 7.8]{DE_Rubin1991} show that
\[
 (\mathcal C_L^{\rm R}(1)_{\Delta\times\Gamma})\otimes\mathbb Q_3
 \longrightarrow
 (\mathcal E_L(1)_{\Delta\times\Gamma})\otimes\mathbb Q_3
                                                               \tag{A.21.4}
\]
is an isomorphism. The same conclusion follows by
taking the invariants and coinvariants of
\(0\to\mathcal C_L^{\rm R}\to\mathcal E_L
\to\mathcal E_L/\mathcal C_L^{\rm R}\to0\): both
end terms from the quotient are finite. Consequently
Rubin units have nonzero image under (A.21.2).

We transfer this image from the cyclotomic tower to
\(k(3^\infty)\). Write \(\mathcal C^{\rm R}\) for
its two-variable elliptic-unit module. The kernel and
cokernel of the norm map
\[
 (\mathcal C^{\rm R})_{\operatorname{Gal}(k(3^\infty)/L_\infty)}
                   \longrightarrow\mathcal C_L^{\rm R}
\]
are annihilated by the product of the annihilator of
\(\mathbb Z_3(1)\) and the augmentation ideals of the
two local decomposition groups, by
\cite[Theorem 7.7(ii)]{DE_Rubin1991}. After twist \((1)\)
and coinvariants, a displayed annihilator evaluates as
\[
 (\chi_3(\sigma)^{-1}-\chi_3(\sigma))
 (\chi_3(\tau_1)^{-1}-1)
 (\chi_3(\tau_2)^{-1}-1).                              \tag{A.21.5}
\]
The cyclotomic character maps each local decomposition
group onto an open subgroup of \(\mathbb Z_3^\times\).
Choose all three character values equal to a common
\(u=4^a\), \(a\ge1\), lying in these images. Then
(A.21.5) is \((u^{-1}-u)(u^{-1}-1)^2\ne0\).
For \(u=4\) its value is \(-135/64\).
Localization at the evaluation character therefore
annihilates the kernel and cokernel before taking
coinvariants, so no nonexact tensor step is being
silently used. The rational norm map on the required
specializations is an isomorphism. Norm compatibility
of Kummer corestriction proves that the two-variable
Rubin units also have nonzero rational Soul\'e image.

Now compare the units actually defining the geometric
character. Let \(\mathcal C^{\theta}\) be formed from
the theta values of \cite[\S2.2]{DE_Ishii2024}.
At every finite level, the definitions give
\[
 12 C^{\rm R}\ \subset\ C^{\theta}\ \subset\ C^{\rm R}
                                                               \tag{A.21.6}
\]
in additive notation, after adjoining the same roots
of unity. A Rubin generator is an appropriate twelfth
root of the corresponding theta-unit generator;
raising it to twelve gives that generator, and theta
units themselves belong to Rubin's group. These
inclusions commute with norms and use the same
integer twelve at every level. Hence (A.21.6) also
holds for the inverse-limit modules and makes their
rationalizations equal. Roots of unity cause no
additional image: after twist \((1)\) they have
character \(\chi_3^2\), whose specialization at the
trivial character is rationally zero. Thus theta
units have nonzero rational Soul\'e image at three.
This is the step for which the published argument
uses \(p\ge5\); replacing an integral equality by the
uniform inclusions (A.21.6) proves the rational
assertion at three.

Finally the distribution and Kummer identities in
\cite[\S3.1 and \S4.1]{DE_Ishii2024} identify that
image with multiples of \(\kappa_{(2,2)}\). We record
the specialization explicitly. Take
\(\pi=1+\sqrt{-2}\), so \(\pi\overline\pi=3\),
and \(w_k=2\). For any auxiliary ideal \(\mathfrak a\)
prime to six, the primitive theta-unit character
\(\kappa_{(2,2),\mathfrak a}\) satisfies
\[
 \underbrace{2\left(\frac1{1-3}+\frac1{1-3}-1\right)}_{-4}
       \kappa_{(2,2),\mathfrak a}
   =\left(N\mathfrak a-(N\mathfrak a)^{-1}\right)
       \kappa_{(2,2)}.                                  \tag{A.21.7}
\]
The identity follows by separating the two coordinates
of the primitive torsion sum, applying the theta
distribution relation, and summing the geometric
series of ratio three. Its denominators \(1-3\)
are three-adic units, so this calculation is valid at
three. The factor on the right is nonzero for every
nontrivial auxiliary ideal. The theta units in the
two-variable inverse limit are generated by the
primitive theta values and their Galois translates:
the one-coordinate units vanish in this inverse
limit, because norms along the other tower raise
those units to arbitrarily high powers of three.
Taking \(\sigma-1\) simply multiplies their specialized
Soul\'e images by \(\chi_3(\sigma)^{-1}-1\).
Thus a zero \(\kappa_{(2,2)}\) would make the entire
nonzero theta-unit image vanish, contradicting the
preceding paragraph. This proves nonvanishing.
For example \(\mathfrak a=(5)\) has norm twenty-five
and the right factor in (A.21.7) is \(624/25\ne0\).

All arguments just given concern rational nonvanishing;
no assertion about integral surjectivity is required.
The character is a continuous homomorphism to a
characteristic-zero three-adic line on its defining
kernel. If its restriction to a finite-index subgroup
were zero, its image would be finite and hence zero.
This proves the final assertion.
\end{proof}
